% Part: proof-theory
% Chapter: cut-elimination
% Section: ce-largest

\documentclass[../../../include/open-logic-section]{subfiles}

\begin{document}

\olfileid[pt]{pt}{cut}{inv}

\olsection{Remoção dos cortes de maior posto}

A demonstração de \olref[seq][inv]{lem:G3c-invert} mostrou que, se \Log{G3c}
!!{prove}s a conclusão de uma regra lógica (que não seja
\RightR{\lexists} nem \LeftR{\lforall}), também !!{prove}s cada uma das
premissas da regra sem aumentar a !!{height} da
demonstração. Podemos obter um resultado mais forte e mostrar que isso também vale para $\Log{G3c} +
\CutCS$.

Como antes, definimos o \emph{posto do corte} de uma inferência \CutCS{} como
a profundidade de sua !!{formula} de corte.

\begin{lem}\ollabel{lem:inv-G3c-cut}
Se $\Log{G3c} + \CutCS$ !!{prove}s a conclusão de uma regra
lógica~$R$ que não seja \RightR{\lexists} nem \LeftR{\lforall}, usando !!a{proof}~$\pi$, então também !!{prove}s cada premissa
com !!{proof}~$\pi'$ (ou !!{proof}s~$\pi'$, $\pi''$, conforme
a regra tenha uma ou duas premissas). Além disso,
(a)~$\pheight{\pi'}\le\pheight{\pi}$ (e
$\pheight{\pi''}\le\pheight{\pi}$) e (b)~o número de inferências \CutCS{},
bem como seus postos, é o mesmo em~$\pi'$ (e $\pi''$)
que em~$\pi$.
\end{lem}

\begin{proof}
Por inspeção das demonstrações de
\cref{pt:seq:inv:lem:G3c-invert,pt:seq:inv:lem:invert-quant}. Essa
demonstração foi feita por indução sobre~$\pheight{\pi}$. No caso base,
substituímos um axioma por outro axioma. Portanto, as condições (a) e~(b) são
trivialmente satisfeitas. Se a última regra de~$\pi$ era~$R$, suas
sub-!!{proof}s imediatas são as !!{proof}s~$\pi_i$ que queremos, e (a) e~(b)
são ambas satisfeitas. Em todos os outros casos, aplicamos a hipótese de
indução às sub-!!{proof}s imediatas de~$\pi$, isto é, às que
levam à premissa ou às premissas da última inferência~$R$, e depois aplicamos
essa mesma regra $R$ aos resultados. As condições (a) e~(b) valem para
as sub-!!{proof}s imediatas da nova demonstração e, portanto, (a) e~(b) valem
para a demonstração inteira. Resta verificar o caso em que a última
inferência é~\CutCS. Fazemos isso novamente para um único exemplo: suponha que
$\pi$ seja !!a{proof} da conclusão $!B \land !C, \Gamma \Sequent
\Delta$ de \LeftR{\land} e termine em~\CutCS:
\[
\AxiomC{}
\RightLabel{$\pi_1$}
\Deduce$!B \land !C, \Gamma \fCenter \Delta, !A$
\AxiomC{}
\RightLabel{$\pi_2$}
\Deduce$!A, !B \land !C, \Gamma \fCenter \Delta$
\RightLabel{\CutCS}
\BinaryInf$!B \land !C, \Gamma \fCenter \Delta$
\DisplayProof
\]
A hipótese de indução se aplica a $\pi_1$ e~$\pi_2$ (que também são
!!{proof}s de instâncias da conclusão de $\LeftR{\land}$) e
fornece !!{proof}s $\pi_1'$ e~$\pi_2'$ de $!B, !C, \Gamma \fCenter
\Delta, !A$ e $!A, !B, !C, \Gamma \fCenter \Delta$,
respectivamente. Combinamos essas demonstrações usando \CutCS{} para obter~$\pi'$:
\[
\AxiomC{}
\RightLabel{$\pi_1'$}
\Deduce$!B, !C, \Gamma \fCenter \Delta, !A$
\AxiomC{}
\RightLabel{$\pi_2'$}
\Deduce$!A, !B, !C, \Gamma \fCenter \Delta$
\RightLabel{\CutCS}
\BinaryInf$!B, !C, \Gamma \fCenter \Delta$
\DisplayProof
\]
É claro que (a) e~(b) são satisfeitas.
\end{proof}

Observe que isso não funcionaria se usássemos \Cut{} em vez de \CutCS, pois
então a conclusão da última !!{proof} teria duas cópias de
$!B, !C, \Gamma$ no antecedente e duas cópias de $\Delta$ no
consequente. Teríamos de recorrer à admissibilidade da
contração, mas a demonstração de \olref[seq][inv]{prop:G3c-cont-adm} usa
o lema de invertibilidade. Estaríamos, portanto, argumentando em círculo.

Podemos usar \cref{pt:cut:inv:lem:inv-G3c-cut} para dar uma demonstração
alternativa da eliminação do corte. Nessa demonstração, não removemos
sucessivamente as ocorrências mais altas de inferências \CutCS{}, mas sim
ocorrências de inferências \CutCS{} de posto máximo. Isto é, começamos com
!!a{proof}~$\pi$ em $\Log{G3c} + \CutCS$ e removemos um corte em qualquer
posição nela (talvez o substituindo por cortes de posto menor)---e então
continuamos a fazer isso até não restar nenhum corte. A única condição é que,
acima do corte de que tratamos, não ocorram cortes de posto igual ou maior.

\begin{lem}\ollabel{lem:max-cut-red-G3c} Se $\pi$ é !!a{proof} em
$\Log{G3c}+\CutCS$ que termina em uma inferência~\CutCS{} de posto~$n$ e,
quanto ao mais, usa apenas regras de~$\Log{G3c}$ e inferências \CutCS{} de
posto $<n$, então existe uma demonstração~$\pi'$ do mesmo sequente final
em~$\Log{G3c} + \CutCS$ na qual toda inferência \CutCS{} tem posto~$<n$.
\end{lem}

\begin{proof}
A !!{proof}~$\pi$ termina em:
\[
\AxiomC{}
\RightLabel{$\pi_1$}
\Deduce$\Gamma \fCenter \Delta, !A$
\AxiomC{}
\RightLabel{$\pi_2$}
\Deduce$!A, \Gamma \fCenter \Delta$
\RightLabel{\CutCS}
\BinaryInf$\Gamma \fCenter \Delta$
\DisplayProof
\]
Distinguimos casos de acordo com a forma de~$!A$.

Suponha que $!A$ seja atômica. Observe que, pela condição de que todas as
inferências \CutCS{} em~$\pi_1$ e $\pi_2$ devem ter posto $<
\depth{!A} = 0$, não pode haver inferências \CutCS{} em~$\pi_1$
nem em~$\pi_2$. Mostramos que $\Gamma \Sequent \Delta$ tem uma demonstração
sem corte por indução sobre $\pheight{\pi_2}$. Se
$\pheight{\pi_2} = 0$, então $!A, \Gamma \Sequent \Delta$ é um axioma.
Se $!A$ não é principal, alguma $!C$ atômica é !!a{element} tanto
de~$\Gamma$ quanto de $\Delta$, e o próprio $\Gamma \Sequent \Delta$ é
um axioma. Se $!A$ é principal, então $\Delta = \Delta, !A$. Nesse caso,
$\pi_1$ termina em $\Gamma \Sequent \Delta', !A, !A$. Obtemos uma
demonstração sem corte de $\Gamma \Sequent \Delta', !A$ aplicando
\olref[seq][inv]{prop:G3c-cont-adm} (admissibilidade da contração). Se
$\pheight{\pi_2} > 0$, ela termina em uma regra~$R$. Tome uma premissa
dessa regra: ela tem a forma $!A, \Gamma', \Pi \Sequent \Delta',
\Lambda$, em que $\Pi$ e $\Lambda$ são as !!{formula}s ativas de~$R$.
Aplicar \olref[seq][adm]{prop:weak-G3c-adm} a~$\pi_1$ e $\pi_2$ fornece
!!{proof}s de $\Gamma, \Gamma', \Pi \Sequent \Delta, \Delta', \Lambda,
!A$ e $!A, \Gamma, \Gamma', \Pi \Sequent \Delta, \Delta', \Lambda$,
às quais se aplica a hipótese de indução. Isso nos dá !!a{proof} de
$\Gamma, \Gamma', \Pi \Sequent \Delta, \Delta', \Lambda$ para cada
premissa de~$R$. Aplicar $R$ fornece $\Gamma, \Gamma \Sequent \Delta,
\Delta$, e \olref[seq][inv]{prop:G3c-cont-adm} fornece a demonstração
sem corte de $\Gamma \Sequent \Delta$.

Se $!A$ não é atômica, distinguimos casos de acordo com o !!{main
operator} de~$!A$. Em cada caso, aplicamos \olref{lem:inv-G3c-cut} para
obter !!{proof}s das premissas das regras à esquerda e à direita que têm
$!A$ como fórmula principal. Procedemos então como no caso~E da demonstração
de \olref[top]{lem:cut-adm-G3c}.
\begin{enumerate}
  \item $!A \ident \lnot !B$. As !!{proof}s $\pi_1$ e $\pi_2$ terminam
  em $\Gamma \Sequent \Delta, \lnot !B$ e $\lnot !B, \Gamma \Sequent
  \Delta$. Por \olref{lem:inv-G3c-cut}, existem !!{proof}s $\pi_1'$
  e $\pi_2'$ de $!B, \Gamma \Sequent
  \Delta$ e $\Gamma \Sequent \Delta, !B$. A !!{proof}~$\pi'$ é:
  \[
  \AxiomC{}
  \RightLabel{$\pi_2'$}
  \Deduce$\Gamma \fCenter \Delta, !B$
  \AxiomC{}
  \RightLabel{$\pi_1'$}
  \Deduce$!B, \Gamma \fCenter \Delta$
  \RightLabel{\CutCS}
  \BinaryInf$\Gamma \fCenter \Delta$
  \DisplayProof
  \]
  \item $!A \ident (!B \lor !C)$. As !!{proof}s $\pi_1$ e $\pi_2$
  terminam em $\Gamma \Sequent \Delta, !B \lor !C$ e $!B \lor !C, \Gamma
  \Sequent \Delta$. Por \olref{lem:inv-G3c-cut} (inversão para
  \RightR{\lor} aplicada a~$\pi_1$), existe !!a{proof} $\pi_1'$ de
  $\Gamma \Sequent \Delta, !B, !C$. Por \olref{lem:inv-G3c-cut}
  (inversão para \LeftR{\lor} aplicada a~$\pi_2$), existe !!a{proof}
  $\pi_2'$ de $!B, \Gamma \Sequent \Delta$ e !!a{proof} $\pi_2''$ de
  $!C, \Gamma \Sequent \Delta$. Aplicando
  \olref[seq][adm]{prop:weak-G3c-adm} a $\pi_2''$, obtemos também
  !!a{proof}~$\pi_2'''$ de $!C, \Gamma \Sequent \Delta, !B$. A
  !!{proof}~$\pi'$ é:
  \[
  \AxiomC{}
  \RightLabel{$\pi_1'$}
  \Deduce$\Gamma \fCenter \Delta, !B, !C$
  \AxiomC{}
  \RightLabel{$\pi_2'''$}
  \Deduce$!C, \Gamma \fCenter \Delta, !B$
  \RightLabel{\CutCS}
  \BinaryInf$\Gamma \fCenter \Delta, !B$
  \AxiomC{}
  \RightLabel{$\pi_2'$}
  \Deduce$!B, \Gamma \fCenter \Delta$
  \RightLabel{\CutCS}
  \BinaryInf$\Gamma \fCenter \Delta$
  \DisplayProof
  \]
  \item $!A \ident (!B \land !C)$. Exercício.
  \item $!A \ident (!B \lif !C)$. As !!{proof}s $\pi_1$ e $\pi_2$
  terminam em $\Gamma \Sequent \Delta, !B \lif !C$ e $!B \lif !C, \Gamma
  \Sequent \Delta$. Por \olref{lem:inv-G3c-cut} (inversão para
  \RightR{\lif} aplicada a~$\pi_1$), existe !!a{proof} $\pi_1'$ de
  $!B, \Gamma \Sequent \Delta, !C$. Por \olref{lem:inv-G3c-cut}
  (inversão para \LeftR{\lif} aplicada a~$\pi_2$), existe
  !!a{proof} $\pi_2'$ de $\Gamma \Sequent \Delta, !B$ e !!a{proof}
  $\pi_2''$ de $!C, \Gamma \Sequent \Delta$. Aplicando
  \olref[seq][adm]{prop:weak-G3c-adm} a $\pi_2''$, obtemos também
  !!a{proof}~$\pi_2'''$ de $!C, !B, \Gamma \Sequent \Delta$. A
  !!{proof}~$\pi'$ é:
  \[
  \AxiomC{}
  \RightLabel{$\pi_2'$}
  \Deduce$\Gamma \fCenter \Delta, !B$
  \AxiomC{}
  \RightLabel{$\pi_1'$}
  \Deduce$!B, \Gamma \fCenter \Delta, !C$
  \AxiomC{}
  \RightLabel{$\pi_2'''$}
  \Deduce$!C, !B, \Gamma \fCenter \Delta$
  \RightLabel{\CutCS}
  \BinaryInf$!B, \Gamma \fCenter \Delta$
  \RightLabel{\CutCS}
  \BinaryInf$\Gamma \fCenter \Delta$
  \DisplayProof
  \]
  \item $!A \ident \lforall[x][!B(x)]$. As !!{proof}s $\pi_1$ e
  $\pi_2$ terminam em $\Gamma \Sequent \Delta, \lforall[x][!B(x)]$ e
  $\lforall[x][!B(x)], \Gamma \Sequent \Delta$. Por
  \olref{lem:inv-G3c-cut}, existe !!a{proof}~$\pi_1'$ de $\Gamma
  \Sequent \Delta, !B(c)$.

  Considere agora a !!{proof}~$\pi_2$. Ela termina em $\lforall[x][!B(x)],
  \Gamma \Sequent \Delta$. Subindo a partir do sequente final
  de~$\pi_2$, marque toda ocorrência de $\lforall[x][!B(x)]$ à
  esquerda de um sequente que ``conduza à'' ocorrência de
  $\lforall[x][!B(x)]$ no sequente final de~$\pi_2$, isto é:
  (a)~Marque $\lforall[x][!B(x)]$ no sequente final de~$\pi_2$. (b)~Se
  $\lforall[x][!B(x)]$ ocorre no contexto da conclusão de uma
  regra e está marcada, marque uma ocorrência correspondente na premissa
  ou nas premissas da regra. (c)~Se $\lforall[x][!B(x)]$ é principal na
  conclusão de uma regra \LeftR{\lforall} e está marcada, marque as
  ocorrências ativas correspondentes de $\lforall[x][!B(x)]$ e $!B(t)$
  na premissa.
  
  Considere agora todas essas ocorrências marcadas de $\lforall[x][!B(x)]$.
  Nenhuma delas é ativa em uma regra diferente de \LeftR{\lforall} (somente
  !!{formula}s ativas de \LeftR{\lforall} são marcadas). Apague todas as
  ocorrências marcadas de $\lforall[x][!B(x)]$ em $\pi_2$.
  Se o resultado for uma !!{proof} correta, as ocorrências marcadas de
  $\lforall[x][!B(x)]$ nunca foram ativas; todas vieram diretamente de
  axiomas. A !!{proof} termina em~$\Gamma \Sequent \Delta$, e podemos
  substituir $\pi$ por ela. Removemos um corte de posto máximo.

  Caso contrário, o que temos não é uma demonstração correta, pois apagamos
  !!{formula}s $\lforall[x][!B(x)]$ que eram ativas em algumas inferências
  \LeftR{\lforall}. Transformamos essas inferências em
  ``inferências'' da forma
  \[
  \AxiomC{}
  \RightLabel{$\pi_t$}
  \Deduce$!B(t), \Pi \fCenter \Lambda$
  \RightLabel{\LeftR{\lforall}}
  \UnaryInf$\Pi \fCenter \Lambda$
  \DisplayProof
  \]
  Substitua cada ``inferência'' desse tipo que seja mais alta por
  \[
  \AxiomC{}
  \RightLabel{$\Subst{\pi_1'}{t}{c}[\Pi \Sequent \Lambda]$}
  \Deduce$\Gamma, \Pi \fCenter \Delta, \Lambda, !B(t)$
  \AxiomC{}
  \RightLabel{$\pi_t[\Gamma \Sequent \Delta]$}
  \Deduce$!B(t), \Gamma, \Pi \fCenter \Delta, \Lambda$
  \RightLabel{\CutCS}
  \BinaryInf$\Gamma, \Pi \fCenter \Delta, \Lambda$
  \DisplayProof
  \]
  e acrescente $\Gamma$ à esquerda e $\Delta$ à direita de cada sequente
  abaixo dela. Repita até que todas as ``inferências'' \LeftR{\lforall}
  com um $!B(t)$ marcado na premissa tenham sido substituídas por uma
  inferência \CutCS{} sobre algum~$!B(t)$. O sequente final da
  !!{proof} resultante (que agora é uma !!{proof} correta) tem a forma
  $\Gamma, \dots, \Gamma \Sequent \Delta, \dots, \Delta$. Aplique a
  admissibilidade da contração para transformá-la em !!a{proof} de~$\Gamma
  \Sequent \Delta$ e substitua $\pi$ por ela. Substituímos um corte
  sobre $\lforall[x][!B(x)]$ por cortes (possivelmente muitos) sobre
  $!B(t)$, mas todos eles têm posto menor. Quaisquer cortes que já ocorressem
  em $\pi_1$ ou $\pi_2$ permanecem presentes, mas também têm todos posto menor.
\end{enumerate}
\end{proof}

\begin{prob}
Verifique o caso em que $!A \ident (!B \land !C)$ na demonstração de
\olref{lem:inv-G3c-cut}. Isto é, mostre que, se existe
!!a{proof}~$\pi$ que termina em 
\[
\AxiomC{}
\RightLabel{$\pi_1$}
\Deduce$\Gamma \fCenter \Delta, !B \land !C$
\AxiomC{}
\RightLabel{$\pi_2$}
\Deduce$!B \land !C, \Gamma \fCenter \Delta$
\RightLabel{\CutCS}
\BinaryInf$\Gamma \fCenter \Delta$
\DisplayProof
\]
com $\pi_1$ e $\pi_2$ contendo apenas cortes de posto $<\depth{!B
\land !C}$, então existe !!a{proof}~$\pi'$ de $\Gamma \Sequent \Delta$
que contém apenas cortes de posto $<\depth{!B \land !C}$.
\end{prob}

% \begin{prob}
% Verify the case where $!A \ident \lexists[x][!B(x)]$ in the proof of
% \olref{lem:inv-G3c-cut}.
% \end{prob}


\olref[top]{lem:cut-adm-G3c} estabelece que podemos substituir uma
inferência \CutCS{} de maior posto em !!a{proof} por inferências \CutCS{} de
posto menor. Podemos usar novamente esse lema para mostrar que podemos eliminar
qualquer número de inferências \CutCS{} de !!a{proof}, aplicando-o
sucessivamente às sub-!!{proof}s mais altas que terminam em~\CutCS{} de posto máximo.

\begin{thm}\ollabel{thm:cut-elim-G3c-largest}
Toda !!{proof} $\pi$ em $\Log{G3c} + \CutCS$ pode ser transformada em uma demonstração em~\Log{G3c}.
\end{thm}

\begin{proof}
  Primeiro, provamos que todos os cortes de posto máximo em~$\pi$ podem ser
  removidos. Seja $d$ o posto máximo dos cortes que ocorrem em~$\pi$, e
  seja $n$ o número de cortes de posto~$d$. A demonstração é por indução
  sobre~$n$. Se $n=0$, não há nada a provar. Se $n>0$, escolha um corte
  de posto~$d$ que seja mais alto e seja $\pi_1$ a sub-!!{proof} que termina
  nele. Por \olref{lem:max-cut-red-G3c}, existe !!a{proof}~$\pi_1'$ do
  mesmo sequente final, com todos os cortes de posto $<d$. Substitua $\pi_1$
  por $\pi_1'$ em~$\pi$. O resultado é !!a{proof} que contém $n-1$ cortes
  de posto~$d$, e a hipótese de indução se aplica.

  Agora o resultado decorre por indução sobre~$d$. Se $d=0$, todos os cortes
  são atômicos e, pelo resultado anterior, podem ser todos removidos. Se
  $d>0$, o resultado anterior fornece uma demonstração na qual todos os
  cortes de posto $d$ foram substituídos por cortes de posto~$<d$. Como $d$
  era o posto máximo dos cortes em~$\pi$, temos uma demonstração na qual o
  posto máximo dos cortes é $<d$, e a hipótese de indução se aplica.
\end{proof}

\end{document}
